% exp-fips203-mlkem-pke-decrypt-k3 — FIPS 203 Alg.15 ML-KEM-768 PKE Decrypt
% 编译：bash scripts/xelatex-clean.sh examples/incubating/ml-kem/ml-kem-768/exp-fips203-mlkem-pke-decrypt-k3/exp-fips203-mlkem-pke-decrypt-k3-实现方案-customspec.tex
\documentclass[UTF8,a4paper,11pt]{ctexart}
\usepackage{amsmath,amssymb}
\usepackage{booktabs}
\usepackage{geometry}
\usepackage{hyperref}
\usepackage{longtable}
\usepackage{array}
\usepackage{seqsplit}
\geometry{margin=2.2cm}
\newcolumntype{L}[1]{>{\raggedright\arraybackslash}p{#1}}
\newcommand{\apiname}[1]{\texttt{\seqsplit{#1}}}
\newcommand{\dirname}[1]{\texttt{\seqsplit{#1}}}
\hypersetup{colorlinks=true,linkcolor=blue,urlcolor=blue}

\title{exp-fips203-mlkem-pke-decrypt-k3\\FIPS 203 Alg.15 ML-KEM-768 PKE Decrypt 实现方案}
\author{ascendc 工程（ML-KEM-768 W4a / E15）}
\date{2026-07-26}

\begin{document}
\maketitle

\begin{abstract}
本文档描述 \dirname{examples/incubating/ml-kem/ml-kem-768/exp-fips203-mlkem-pke-decrypt-k3/}：
在 \textbf{Atlas A2 / Ascend910B4} 上以 AscendC 实现 FIPS 203 \textbf{Algorithm 15}，即 ML-KEM-768（$k=3$）\textbf{PKE Decrypt}。

\textbf{参数锁定}：几何、I/O 与 launch 全部锁定到
\dirname{docs/specs/fips203-mlkem768-parameter-card.md} §3.2 D15。
\textbf{行为基线}：活跃探针
\dirname{ascendc-tests/ml-kem/ml-kem-768/pass-fix-f203-alg15-pke-decrypt-device-k3/}。
\textbf{生产 I/O}：\texttt{input/} 为 \texttt{dk\_pke.bin}（1152B）、\texttt{c.bin}（1088B）与静态 LUT；
\texttt{output/} 仅 \texttt{m.bin}（32B）。
\textbf{计划 AI Core 数}：1 launch；\texttt{f203\_decrypt\_device\_fused} 为 MIX \texttt{blockDim=1}（1 AIC + 2 AIV）。
\textbf{禁止}：创建 \dirname{examples/stable/ml-kem/ml-kem-768/}；从任何 \dirname{frozen/} 目录复制、改写或移植源码；用 k4 的 \texttt{dk=1536}、\texttt{c=1568}、\texttt{du/dv=11/5} 或 4/8 路假 poly 替代 k3 几何。
\end{abstract}

\tableofcontents
\newpage

\section{工程边界}

\subsection{允许的代码来源与依赖}

\begin{longtable}{@{}L{4.6cm}L{8.6cm}@{}}
\toprule
来源 & 用途 \\
\midrule
\dirname{ascendc-tests/ml-kem/ml-kem-768/pass-fix-f203-alg15-pke-decrypt-device-k3/} & 活跃 D15 k3 行为基线；可复制后自包含适配到本 exp，最终不得保留对探针目录的编译期 include/import \\
\dirname{examples/incubating/ml-kem/ml-kem-1024/exp-fips203-mlkem-pke-decrypt-k4/} & 活跃 k4 incubating 结构参考；仅 retarget 到 k3 锁定几何与字节长 \\
\dirname{library/shared/} & 公共设备侧头、统一整数压缩/解压辅助等共享库 \\
CANN / tikicpulib / CAModel & 编译、CPU 孪生与 SIM 运行环境 \\
\bottomrule
\end{longtable}

\subsection{禁止项}

\begin{itemize}
  \item 禁止从 \dirname{ascendc-tests/frozen/} 或 \dirname{examples/frozen/} 拿源码、customspec 或路线。
  \item 禁止将 \(u'\)、\(v'\)、\(\hat s\)、\(\hat u\)、\(\hat w\)、\(w\) 等中间态作为默认 I/O 落盘；debug dump 只能显式开关启用。
  \item 禁止把 k4 的 \texttt{dk=1536}、\texttt{c1=1408/c2=160}、\texttt{du=11/dv=5}、4/8 路 pad 当作默认。
  \item 禁止新建或写入 stable-768；本目录只作为 incubating 预研实现。
\end{itemize}

\section{数学语义（Alg.15，$k=3$，$n=256$，$q=3329$）}

\subsection{输入输出}

\begin{longtable}{@{}L{3.3cm}L{10.2cm}@{}}
\toprule
张量 / 文件 & 契约 \\
\midrule
\texttt{input/dk\_pke.bin} & 1152B；\(\hat s\) 的 \(\mathrm{ByteEncode}_{12}\)，即 \(3\times384\)B \\
\texttt{input/c.bin} & 1088B；\(c=c_1\Vert c_2\)，其中 \(c_1=960\)B（\(d_u=10\)，3 poly），\(c_2=128\)B（\(d_v=4\)，1 poly） \\
\texttt{input/lut\_even\_stacked.bin}、\texttt{lut\_odd\_stacked.bin} & NTT/INTT Stage2 LUT；由本目录 golden 脚本生成；不是用户可见交付输入 \\
\texttt{output/m.bin} & 32B，\(\mathrm{ByteEncode}_{1}(\mathrm{Compress}_{1}(w))\)，为唯一产物 \\
\bottomrule
\end{longtable}

\subsection{阶段公式}

\begin{enumerate}
  \item 行 1--2：拆分 \(c_1=c[0:960]\)，\(c_2=c[960:1088]\)。
  \item 行 3：\(u'\leftarrow\mathrm{Decompress}_{10}(\mathrm{ByteDecode}_{10}(c_1))\)，得到 3 个 poly。
  \item 行 4：\(v'\leftarrow\mathrm{Decompress}_{4}(\mathrm{ByteDecode}_{4}(c_2))\)，得到 1 个 poly。
  \item 行 5：\(\hat s\leftarrow\mathrm{ByteDecode}_{12}(\texttt{dk\_pke})\)，得到 3 个 NTT 域 poly。
  \item 行 6：\(\hat u\leftarrow\mathrm{NTT}(u')\)；\(\hat w\leftarrow\sum_{j=0}^{2}\mathrm{MultiplyNTTs}(\hat s[j],\hat u[j])\)；再对 \(\hat w\) pad 到 k3 前缀后执行 INTT，得到 \(w_t\)。
  \item 行 7：\(w=(v'-w_t)\bmod q\)，\(m=\mathrm{ByteEncode}_{1}(\mathrm{Compress}_{1}(w))\)。
\end{enumerate}

\textbf{尾段方向锁定}：Decrypt 行 7 是 \(\mathrm{ByteEncode}_{1}(\mathrm{Compress}_{1}(w))\)，不是 Encrypt 行 20 的 \(\mathrm{Decompress}_{1}(\mathrm{ByteDecode}_{1}(m))\)。

\section{AscendC 实现规划}

\subsection{Launch 与 AI Core 规模}

\begin{longtable}{@{}L{2.5cm}L{3.0cm}L{7.6cm}@{}}
\toprule
Launch & 核类型 / blockDim & 数据划分与理由 \\
\midrule
1 fused & MIX，\texttt{blockDim=1} & 1 AIC + 2 AIV；AIV0 负责 unpack、\(\hat s\cdot\hat u\) 与尾 pack 的关键串行段；NTT/INTT 使用真实 k3 分片，AIV0 处理 poly0/1，AIV1 处理 poly2。选择单 MIX launch 是参数卡 §3.2 D15 锁定值，继承 D15 探针已验 GATE 4/8 与 softSyncGm。 \\
\bottomrule
\end{longtable}

\subsection{GM 几何}

\begin{longtable}{@{}L{3.7cm}L{9.5cm}@{}}
\toprule
对象 & 锁定几何 \\
\midrule
\(u'\) & \texttt{[3,256] int32}；由 c1 d10 unpack + Decompress10 得到 \\
\(v'\) & \texttt{[1,256] int32}；由 c2 d4 unpack + Decompress4 得到 \\
\(\hat s\) & \texttt{[3,256] int32}；由 dk\_pke ByteDecode12 得到 \\
\(\hat u\) NTT & 真 polyvec3；AIV 2+1，禁止读写第 4 个假 poly \\
\(\hat w\) / INTT & \(\hat w\) 写入 k3 前缀，后续 INTT 只消费第 0 个时域 poly；Stage123 内不引入 k4 batch \\
Tail pack & Compress1 产 256 bit；ByteEncode1 LSB-first pack 为 32B \\
\bottomrule
\end{longtable}

\subsection{核内同步}

单 kernel 内部沿用活跃 D15 k3 已验 FSM：NTT 使用 CrossCore flag 1/2/3；su\_dot 后通过 softSyncGm 汇合双 AIV；段间用 GATE 4\(\to\)8 分隔 NTT 与 INTT；INTT 使用 flag 1/3，避免复用 NTT 的 flag 2。
CPU 仅作为逻辑孪生；SIM 是同步与搬运验收门槛。

\section{AscendC API 表}

本实现不引入 E15 独有的新 AscendC API；下表 API 均已有索引记录，使用约束按
\dirname{library/documents/CANN-AscendC算子开发接口参考-查阅索引.md} 执行。

\begin{longtable}{@{}L{3.2cm}L{2.7cm}L{2.5cm}L{2.2cm}L{3.0cm}@{}}
\toprule
API 名 & 分类 / 文档 & Atlas A2 & 本用例 dtype & 备注 \\
\midrule
\apiname{GetBlockIdx} / \apiname{GetSubBlockIdx} & 系统变量 / 资源 & √ & 标量索引 & 区分 AIV/AIC 子核与串行 owner；禁止把 CPU 打印误认为多核性能 \\
\apiname{GlobalTensor} / \apiname{LocalTensor} & 基础结构 & √ & \texttt{uint8/int8/int32} & GM/UB 张量契约必须与本 spec 几何一致 \\
\apiname{DataCopy} / \apiname{Duplicate} & Memory / 初始化 & √ & \texttt{uint8/int8/int32} & pad \(\hat w\)、清 workspace、GM/UB 交接；禁止标量写 GM 再 MTE 读 \\
\apiname{PipeBarrier} & Pipe 同步 & √ & — & MTE/Vector/Cube 阶段交界显式同步 \\
\apiname{CrossCoreSetFlag} / \apiname{CrossCoreWaitFlag} & MIX 同步 & √ & — & NTT/INTT AIV/AIC 通过 GM 交接时使用；CPU 过不能删 \\
\apiname{Mmad} / \apiname{Fixpipe} & 矩阵计算 & √ & \texttt{int8*int8->int32} & Stage2 MMAD；逻辑 k=3，硬件 M 对齐垫不是假 poly \\
\apiname{ShiftRight} & 矢量算术 & √ & \texttt{int32} & Decompress/Compress 与 Barrett 辅助右移 \\
\apiname{Muls} / \apiname{Add} / \apiname{Sub} & 矢量算术 & √ & \texttt{int32/int16} & Decompress、\(v'-w\)、Compress1 辅助 \\
\apiname{Cast} & 类型转换 & √（受限） & \texttt{int32/int16/half/int8} & 遵守索引中 A2 Cast 链约束；不得假设存在 \texttt{int32->int8} 单跳 \\
\apiname{GetValue} / \apiname{SetValue} & LocalTensor 标量访问 & √ & \texttt{uint8/int32} & 仅用于 ByteEncode1 bit pack 等标量缺口 \\
\bottomrule
\end{longtable}

\subsection{评估过但未采用}

\begin{longtable}{@{}L{3.8cm}L{9.2cm}@{}}
\toprule
API / 路线 & 不采用原因 \\
\midrule
真向量 ByteEncode1 bit pack & d=1 跨字节 bit 流，已有标量 pack 更简单且 D15 已验；本轮不把尾段作为性能变量 \\
高阶 \apiname{Matmul<>} & 本路线继承 D15 手写 MMAD + 平面 mat\_c 契约；高阶 Matmul 输出布局与 MIX 交接不作为本轮变量 \\
\apiname{Gather} 用于 NTT S1--S3 & 参数卡与 ML-KEM NTT 定稿禁止；NTT 三段保持平面 mat\_c + 连续 poly 批处理 \\
\bottomrule
\end{longtable}

\section{数学步骤覆盖率}

\begin{longtable}{@{}L{4.2cm}L{4.0cm}L{2.0cm}L{3.0cm}@{}}
\toprule
数学步骤 & 实现方式 / API & 覆盖 & 缺口说明 \\
\midrule
ByteDecode/Decompress d10/d4 & 标量 unpack + 统一整数向量 Decompress & 部分向量 & bit unpack 标量缺口沿用 D15 已验实现 \\
ByteDecode12(\(\hat s\)) & 标量 unpack 到 GM/UB & 部分向量 & 与 D15 基线一致；不引入新向量 unpack \\
NTT(\(u'\)) 与 \(\hat s\cdot\hat u\) & MIX AIV/AIC + \apiname{Mmad} + Alg.11 & 100\% 设备 & NTT 内无 \apiname{Gather}；su\_dot 串行 owner 明确 \\
INTT(\(\hat w\)) & MIX AIV/AIC + Stage3 merge & 100\% 设备 & pad 到 k3 前缀，不读写假第 4 poly \\
Compress1 + ByteEncode1 & Compress1 向量 Barrett + 标量 bit pack & 部分向量 & ByteEncode1 bit pack 标量缺口沿用 D15 已验实现 \\
\bottomrule
\end{longtable}

\section{验证计划}

\subsection{默认验收}

\begin{verbatim}
cd examples/incubating/ml-kem/ml-kem-768/exp-fips203-mlkem-pke-decrypt-k3
bash run.sh -r cpu -v Ascend910B4
SIM_DIRECT=1 bash run.sh -r sim -v Ascend910B4
\end{verbatim}

期望：\texttt{m.bin} 为 32B；与 golden 对拍 \texttt{[verify] PASS max=0 (32 bytes)}。
SIM 通过后记录 \texttt{Total tick} 到 \texttt{qa/active\_sim\_regress\_summary.md}，并检查用例根目录无 stray \texttt{core*.dump}、\texttt{profile\_*log*.toml}。

\subsection{可选对照}

后续可补 D14\(\to\)D15 PKE round-trip，但 E15 本轮门禁为 CPU + \texttt{SIM\_DIRECT=1} sim。
liboqs 仅作为黑盒 I/O oracle；禁止把 liboqs 源码实现同构移植进 AscendC。

\section{产出物}

\begin{itemize}
  \item \dirname{exp-fips203-mlkem-pke-decrypt-k3-实现方案-customspec.tex}
  \item \dirname{exp-fips203-mlkem-pke-decrypt-k3-实现方案-customspec.pdf}
  \item 后续实现文件须与本 customspec 同目录自包含，且自研 Python/C/AscendC 写详细中文注释。
\end{itemize}

\end{document}
